\section{Local partial resolutions and simultaneous families}\label{sec03:partial}

We construct the local modifications used to compare deformation spaces.
Write $C_d$ for the anticanonical cone over $\dP_d$, and write $Q$ for the
anticanonical cone over $\PP^1\times\PP^1$.  The latter is different from
the ordinary double point: its polarization is $\cO(2,2)$, rather than
$\cO(1,1)$.  All morphisms below are proper over a sufficiently small
representative of the cone germ.  The constructions are algebraic, so
they also define modifications of analytic germs.

Fix a nowhere vanishing volume form on the regular locus of the global
threefold, and use its pullback on every crepant modification.  All
identifications of tangent classes with homology classes in this section
use this single form and Lemma~\ref{lem:localhodge}.  In particular,
coordinates introduced in polynomial models are initially algebraic
coordinates.  Their conversion into Hodge-normalized coordinates will
be stated separately.

\subsection{The degree-five discriminant and boundary periods}\label{sec03:degree5}

The degree-five cone has a smooth local deformation base, so it will
remain unchanged on the global partial resolution.  Its discriminant
requires five separate nonvanishing conditions.

\begin{proposition}\label{prop03:degree5family}
The full semiuniversal deformation base of $C_5$ is a smooth
four-dimensional germ.  It admits linear coordinates in which the
reduced discriminant is the union of five distinct hyperplanes.
Every noncentral singular fibre has only ordinary double points,
one for each discriminant hyperplane containing its parameter.
\end{proposition}

\begin{proof}
Let $V=\CC^5$, $A=\bigwedge^2V$, and
\[
 D=\{p\in A:p\wedge p=0\}=\Cone(\operatorname{Gr}(2,V)).
\]
Choose a general six-dimensional subspace $W\subset A$.
The surface $E=\operatorname{Gr}(2,V)\cap\PP(W)$ is smooth by
Bertini.  Adjunction gives $K_E=-H$, and its degree is five.
All smooth degree-five del Pezzo surfaces over $\CC$ are isomorphic,
so its anticanonical cone is the required germ.  Put
$U=W^\perp\subset A^*$ and $B=A/W=U^*$.  The quotient map
\begin{equation}\label{eq03:degree5family}
 L:D\longrightarrow B
\end{equation}
varies the four cutting linear forms by constants and has central
fibre $C_5$.

The Grassmannian cone is Cohen--Macaulay of dimension seven.
For example, its five quadratic Pfaffians have the graded resolution
\[
 0\longrightarrow S(-5)\longrightarrow S(-3)^5
 \longrightarrow S(-2)^5\longrightarrow S
 \longrightarrow\cO_D\longrightarrow0,
 \qquad S=\operatorname{Sym}(A^*).
\]
The four components of $L$ cut a three-dimensional central fibre
and hence form a regular sequence at the vertex.  Miracle flatness
over the smooth four-dimensional base proves flatness there.
Homogeneity then proves flatness everywhere: scaling transports
each point into this neighborhood and acts compatibly on $B$.

We check the full Kodaira--Spencer map.  A constant vector $c\in A$
gives the perturbation $p\wedge c$ of the Pfaffian equations on $W$,
up to the common factor two.  If this abstract first-order class is
zero, the cotangent presentation of the embedded germ shows that
its homogeneous degree-$-1$ part comes from a constant translation
$w\in W$.  The quadratic ideal has no linear elements, so
\[
 p\wedge(c-w)=0\quad\hbox{for every }p\in W.
\]
Write $a=c-w$.  If $a$ has rank four, the map
$a\wedge-:\bigwedge^2V\to\bigwedge^4V$ has rank five; its
kernel cannot contain the six-dimensional space $W$.
If $a$ has rank two with support $P$, its kernel is
$P\wedge V$, of dimension seven.  The decomposable bivectors in
$\PP(P\wedge V)$ form the four-dimensional Schubert variety of
two-planes meeting $P$: choose a line in $P$ and then a second
direction modulo that line.  A hyperplane $\PP(W)$ in this
$\PP^6$ would meet the Schubert variety in dimension at least
three, contradicting $\dim E=2$.  Thus $a=0$.
The Kodaira--Spencer map $B\to T^1_{C_5}$ is injective and is
an isomorphism because $\dim T^1_{C_5}=4$
\cite[\S5.5]{Gross97rev}.  This also shows that all of $T^1_{C_5}$
has degree $-1$.

Embed a semiuniversal base minimally with ring
$\CC\{z_1,\ldots,z_4\}/I$, $I\subset(z_1,\ldots,z_4)^2$.
The classifying map of~\eqref{eq03:degree5family} pulls the $z_i$
back to four functions with independent linear terms.  They are
analytic coordinates on $B$.  Hence the map from the ambient
convergent power-series ring is injective, forcing $I=0$.
The family is therefore semiuniversal for the full germ, including
its scheme structure; local prorepresentability is not required.

For a nonzero decomposable $p$ with support $P_p$, one has
$T_pD=P_p\wedge V$.  An alternating form $\alpha\in A^*$
annihilates this tangent space precisely when
$P_p\subset\ker\alpha$.  Such a nonzero form has rank two.
Thus the dual Grassmannian is the rank-two locus in $\PP(A^*)$.
It meets $\PP(U)$ transversely in five points
$[\alpha_1],\ldots,[\alpha_5]$.  Indeed, for a rank-two
$\alpha\in U$, failure of transversality is equivalent to
\[
 W\cap\bigwedge^2\ker\alpha\ne0.
\]
Every nonzero bivector in this intersection is decomposable and
would be a nonvertex singular point of $C_5$, contrary to the
smoothness of $E$.  The same argument excludes a positive-dimensional
intersection.  Its length is the degree five of the Grassmannian.

Set $K_i=\ker\alpha_i$ and $H_i=\ker(\alpha_i:B\to\CC)$.
The critical locus of $L$ away from the vertex is the union of
the punctured three-dimensional vector spaces $\bigwedge^2K_i$.
The maps
\[
 L:\bigwedge^2K_i\longrightarrow H_i
\]
are isomorphisms: their kernels are the intersections just excluded,
and both dimensions are three.  Thus a point $b\in H_i\setminus0$
has one associated critical point $p_i(b)$, linear in $b$.
Two such points cannot coincide for $b\ne0$.  Otherwise every
form in the pencil spanned by $\alpha_i,\alpha_j$ would vanish
on the same two-plane and hence have rank at most two, giving a
line in the finite dual intersection.  The complete reduced
discriminant is consequently $\bigcup_iH_i$.

Finally, near a nonzero $p_i(b)$ choose a Grassmannian chart in
which $\alpha_i=e_4^*\wedge e_5^*$ and the varying plane is
spanned by
\[
 e_1+x_3e_3+x_4e_4+x_5e_5,\qquad
 e_2+y_3e_3+y_4e_4+y_5e_5.
\]
With the nonzero cone coordinate $r$, the function $\alpha_i$
is $r(x_4y_5-x_5y_4)$.  Three other components of $L$ restrict
to coordinates on its critical locus, by the preceding isomorphism.
Eliminate them.  The remaining equation has nondegenerate quadratic
part in $x_4,x_5,y_4,y_5$, so the holomorphic Morse lemma gives
an ordinary double point with transverse smoothing parameter
$\alpha_i(b)$, up to a unit.
\end{proof}

\begin{lemma}\label{lem03:degree5marking}
Under the local Hodge comparison and the conic-pencil coordinates
\eqref{eq:degree5coordinates}, the discriminant of
Proposition~\ref{prop03:degree5family} is
\begin{equation}\label{eq03:degree5discriminant}
 \prod_{i=1}^5\lambda_i=0,
 \qquad \sum_{i=1}^5\lambda_i=0.
\end{equation}
\end{lemma}

\begin{proof}
Permuting the four blown-up points gives the $S_4$ action on the
first four conic pencils.  The quadratic Cremona involution at
the first three points lifts to an automorphism of $E$ and
interchanges the fourth and fifth pencils, fixing the first three.
These automorphisms generate $S_5$.  On
$\operatorname{Tot}(K_E)$ their canonical lifts preserve the
tautological two-form $\eta$ and the volume form $d\eta$.
The local Hodge comparison is therefore equivariant and realizes
$T^1_{C_5}$ as the standard sum-zero representation on the five
coordinates $\lambda_i$.

The tangent cone of the discriminant is intrinsic to the
deformation functor.  An automorphism of the central germ induces
a classifying map on a semiuniversal base; its invertible derivative
is the natural action on $T^1$.  It preserves the discriminant
because it preserves the isomorphism classes of the fibres.
Hence $S_5$ permutes the five hyperplanes in
Proposition~\ref{prop03:degree5family}.  This permutation action
is faithful.  A nontrivial kernel would contain $A_5$ and make
each normal line invariant under $A_5$.  As $A_5$ has no
nontrivial characters and no invariants in the standard
four-dimensional representation, this is impossible.

The image of this faithful action has order $120$, so it is the
full symmetric group on the five hyperplanes.  The stabilizer
of one hyperplane has order $24$.  In its action
on the five conic pencils it must fix a pencil: without a fixed
point its orbits would have sizes $2+3$, giving order at most
$2!3!=12$, while a size-five orbit cannot divide $24$.
It is thus the $S_4$ stabilizing one pencil.  The restriction
of the standard representation to this subgroup is the sum
of a trivial line and its three-dimensional standard representation.
The unique invariant normal line is generated by the corresponding
$\lambda_i$.  This identifies all five hyperplanes.

For a global volume form $g\,d\eta$, multiplication by the
analytic unit $g$ acts on $T^1_{C_5}$ by $g(0)$: its degree-$-1$
concentration gives $\mathfrak m_0T^1_{C_5}=0$.
Thus the fixed global normalization changes the comparison by
one nonzero scalar and preserves the five hyperplanes.
\end{proof}

\begin{lemma}\label{lem03:degree5milnor}
For a smooth Milnor fibre $M$ of $C_5$ with boundary link $L$,
\[
 H_2(M;\QQ)=0,\qquad \dim H_3(M;\QQ)=4.
\]
Both maps
\begin{equation}\label{eq03:degree5boundary}
 H_3(L;\QQ)\longrightarrow H_3(M;\QQ),\qquad
 H_3(M,L;\QQ)\longrightarrow H_2(L;\QQ)
\end{equation}
are isomorphisms.
\end{lemma}

\begin{proof}
For $b$ outside the discriminant the affine fibre compactifies as
\[
 V_b=\operatorname{Gr}(2,V)\cap\PP(L^{-1}(\CC b)),
 \qquad X_b=V_b\setminus E.
\]
This is a smooth three-hyperplane section of the Grassmannian.
It is smooth at $E$ because the four-hyperplane section $E$ is
smooth, and away from $E$ by Proposition~\ref{prop03:degree5family}.
Near the compact smooth divisor $E$, the squared affine norm is
$F/|s|^2$, where $s$ is a local equation for $E$ and $F$ is
smooth and positive.  Its normal derivative does not vanish
sufficiently close to $E$.  Thus every sufficiently large
Euclidean truncation of $X_b$ removes only a collar at infinity
and has its homotopy type.  Scaling by $\varepsilon$ identifies
$X_b\cap\{\|p\|\le R\}$ with
$X_{\varepsilon b}\cap\{\|p\|\le\varepsilon R\}$.
Fix the latter radius and take $R$ large.  This gives a small
Milnor fibre; transport in the smooth local family gives the
assertion for every nearby smooth parameter.

Lefschetz gives $b_1(V_b)=0$ and $b_2(V_b)=1$.
For completeness the middle Betti number follows from the Chern
calculation.  Write $H=c_1(\mathcal S^*)$ and
$P=c_2(\mathcal S^*)$ for the tautological bundle on the
Grassmannian.  Its tangent bundle is $\mathcal S^*\otimes\mathcal Q$,
so
\[
 c_1=5H,\qquad c_2=11H^2+P,\qquad c_3=15H^3.
\]
Dividing by $(1+H)^3$ gives $c_3(TV_b)=2H^3-3HP$.
Pieri's rule gives
$\int_{V_b}H^3=5$ and $\int_{V_b}HP=2$; hence
$\chi(V_b)=4$ and $b_3(V_b)=0$.
The Gysin sequence for $E\subset V_b$ gives
$H^2(X_b;\QQ)=0$ and
\[
 H^3(X_b;\QQ)\cong
 \ker\{H^2(E;\QQ)\longrightarrow H^4(V_b;\QQ)\},
\]
of dimension four.  The link circle-bundle sequence likewise
gives $b_2(L)=b_3(L)=4$.  Lefschetz duality and the pair sequence
make the second map in~\eqref{eq03:degree5boundary} a surjection
between four-dimensional spaces.  It is an isomorphism.  Exactness
then makes the first map surjective and hence an isomorphism as well.
\end{proof}

\begin{lemma}[degree-five boundary periods]\label{lem03:degree5periods}
There are rational classes $\beta_1,\ldots,\beta_5\in H_3(L;\QQ)$
whose pairings with the local Hodge classes are nonzero multiples
of the five functionals $\lambda_i$.
For any convergent smoothing arc of $C_5$, equipped with a relative
volume form generating its dualizing sheaf, the periods of the
images of all five classes in the smooth fibre are nonzero for
every sufficiently small noncentral parameter.
\end{lemma}

\begin{proof}
Use the semiuniversal family~\eqref{eq03:degree5family}.
The boundary family of a small representative is a smooth
proper submersion over the entire small base $B$, including $0$;
its homology local system is therefore trivial.
By~\eqref{eq03:degree5boundary}, every absolute Milnor cycle
has a unique boundary class.  Near a general point of the
discriminant hyperplane $H_i$, choose the class $\beta_i$
whose image is the nodal vanishing sphere, and transport it in
the boundary family over $B$.

The Pfaffian resolution above shows that the Grassmannian cone
has a homogeneous dualizing generator of weight five.  Dividing
by the four base differentials gives a nowhere vanishing
relative generator $\Omega_b^0$ of weight one.
Integration on transported compact boundary cycles gives
holomorphic functions on the whole base,
\[
 P_i(b)=\int_{\beta_{i,b}}\Omega_b^0.
\]
Holomorphic dependence uses only the smooth locus near the
boundary, so it holds also at parameters with a singular interior.
Scaling changes the boundary radius; compare the two boundaries
through their annular collars, whose central radial isotopy fixes
the homology marking.  The connected scaling action then preserves
transported classes and gives $P_i(cb)=cP_i(b)$.  Thus $P_i$ is linear.
Near a general point of $H_i$, it equals the nodal vanishing
period and has a simple zero along $H_i$.  Consequently
\begin{equation}\label{eq03:degree5linearperiod}
 P_i(b)=c_i\alpha_i(b),\qquad c_i\ne0.
\end{equation}

We verify that its derivative is the claimed Hodge functional.
Choose first-order product trivializations on the punctured
central germ, with Kodaira--Spencer cocycle $\theta_{jk}$.
Differentiating the relative three-form gives
\[
 \dot\Omega_k-\dot\Omega_j
       =d(\iota_{\theta_{jk}}\Omega_0^0).
\]
The \v{C}ech--de Rham representative for the derivative of a
period therefore has overlap term
$\iota_{\theta_{jk}}\Omega_0^0$, precisely the contraction
cocycle defining Lemma~\ref{lem:localhodge}.
In the logarithmic de Rham model of the circle bundle, this is
the natural comparison underlying \cite[Theorem~2.1(vi)]{FL22a}.
The link has $H^3(L;\CC)$ of pure Hodge type $(2,2)$, so its
$F^3$ is zero and this comparison determines the entire class.
Thus $dP_i|_0$ pairs $\beta_i$ with the local Hodge class,
up to the one common convention for Gysin maps and orientation.
There is no ambiguity from the chosen extension of the volume
form: a difference with the same overlap cocycle is a global
form $a\Omega_0^0$.  Normality extends $a$ across the vertex,
and contraction with the Euler vector field makes this form
exact, since all its homogeneous weights are strictly positive.
The same identity applies to a convergent series after shrinking.
Lemma~\ref{lem03:degree5marking} and
\eqref{eq03:degree5linearperiod} identify its kernel with
$\lambda_i=0$.

Finally let $b=b(t)$ be an arbitrary smoothing arc.  Its actual
relative volume form is $g(p,t)\Omega^0$, where $g$ is an
analytic unit on the pulled-back family.  Lift $g$ to a holomorphic
function of the ambient coordinates $(p,t)$; after shrinking it
is a unit and defines the same multiplier over $B\times\Delta$.
The boundary periods are holomorphic functions $P_i^g(b,t)$.
They vanish on $H_i\times\Delta$, by the nodal vanishing-cycle
calculation at its general points and analytic continuation.
Hence
\[
 P_i^g(b,t)=\alpha_i(b)u_i(b,t).
\]
To compute its linear term without expanding $g$ term by term,
fix a smooth direction $b_*$ and a compact cycle $\delta$
representing $\beta_i$ in $X_{b_*}$.  Scaling this cycle gives
\[
 P_i^g(sb_*,t)
 =s\int_\delta g(sp,t)\Omega_{b_*}^0
 =s\,g(0,t)c_i\alpha_i(b_*)+O(s^2).
\]
The estimate is uniform on the fixed compact cycle.  Smooth
directions form a dense open set, so
$d_bP_i^g(0,t)=g(0,t)c_i\alpha_i$.  Thus
$u_i(0,0)=g(0,0)c_i\ne0$.  Restricting to $b(t)$ proves
nonvanishing, whatever its contact order with $H_i$.
\end{proof}


\subsection{The quadric surface cone}

\begin{lemma}\label{lem03:quadric}
The full semiuniversal deformation base of $Q$ is a smooth
one-dimensional germ.  Its general fibre is smooth.  On the crepant
resolution $\operatorname{Tot}(K_{\PP^1\times\PP^1})$, the local Hodge
class of its tangent line is the difference of the two ruling classes.
\end{lemma}

\begin{proof}
The height-one polygon of $Q$, translated to its centre, has vertices
\[
 (1,0),\quad(0,1),\quad(-1,0),\quad(0,-1).
\]
Its oriented edges are $u,v,-u,-v$,
where $u,v$ are linearly independent over $\QQ$.  The polygon has
primitive edges, and the associated threefold singularity is isolated.
Altmann's equations for the versal base are
\[
 \sum_{i=1}^4 t_i^k d_i=0\quad(k\geq1),
 \qquad(d_1,d_2,d_3,d_4)=(u,v,-u,-v),
\]
modulo common translation of the $t_i$.  The equations with $k=1$
are $t_1=t_3$ and $t_2=t_4$.  All higher equations belong to the ideal
generated by these two linear equations.  Consequently the quotient
base is a smooth affine line, scheme-theoretically.  The full versality
theorem, including the obstruction calculation in
\cite[Sections~6--7]{Altmann97}, applies here; the assertion concerns
the entire base, including its possible nilpotent structure.

A smoothing can also be seen directly as
\begin{equation}\label{eq03:quadfamily}
 \{xy-zw=s\}/\boldsymbol\mu_2\longrightarrow\AAA^1_s,
\end{equation}
where the involution negates all four coordinates and fixes $s$.
Taking invariants is exact in characteristic zero, so this quotient
family is flat.  Its central coordinate ring is the second Veronese
subring of the ordinary quadric cone, hence is the section ring of
$\cO_{\PP^1\times\PP^1}(2,2)$.  For $s\ne0$ the quadric is smooth
and the involution acts freely.  Thus the family is a smoothing.

This family is semiuniversal, including its parameter normalization.
Translate the displayed polygon by $(1,0)$ to obtain
$[0,(1,1)]+[0,(1,-1)]$.  Its Cayley cone has rays
\[
 (0,0,1,0),\quad(1,1,1,0),\quad
 (0,0,0,1),\quad(1,-1,0,1).
\]
They span an index-two sublattice of $\ZZ^4$, and half their sum
is integral.  The corresponding toric fourfold is therefore
$\AAA^4/\boldsymbol\mu_2$, with the involution negating all four
coordinates.  The two projection characters pull back to the
products of the first two and the last two coordinates.
Altmann's component family \cite[\S8.1]{Altmann97} is their
difference, precisely~\eqref{eq03:quadfamily} after relabeling.
Its Kodaira--Spencer class is the nonzero difference of the two
segment-dilation parameters.  Since the full base is a smooth line,
every sufficiently small analytic deformation of $Q$ is locally a
pullback of this family.

Finally, the proof of Lemma~\ref{lem:localhodge} applies without change
to $E=\PP^1\times\PP^1$: the same toric vanishing proves the local
cohomology comparison.  The circle-bundle sequence identifies its link
homology with
\[
 K_E^\perp\subset H_2(E;\CC)
   =\CC f_1\oplus\CC f_2,
 \qquad K_E=-2f_1-2f_2.
\]
This perpendicular space is $\CC(f_1-f_2)$ and has dimension one,
which equals the tangent dimension just calculated.  The comparison
therefore identifies the tangent line with this ruling difference.
\end{proof}

\begin{lemma}\label{lem03:quadricmilnor}
Let $M$ be a smooth Milnor fibre of $Q$, with boundary $L$.
Then $M$ is homotopy equivalent to $\mathbb{RP}^3$ and
$\pi_1(L)=\ZZ/2$.  In particular,
\[
 H_2(M;\QQ)=0,\qquad H_3(M;\QQ)=\QQ,\qquad H^1(L;\ZZ)=0.
\]
The boundary map
$H_3(M,L;\QQ)\to H_2(L;\QQ)$ is an isomorphism; under the
circle-bundle identification its image is $\QQ(f_1-f_2)$.
\end{lemma}

\begin{proof}
After a linear change of coordinates the double cover of a
noncentral fibre of~\eqref{eq03:quadfamily} is
$\sum_{i=1}^4z_i^2=s$.  Its Milnor fibre retracts equivariantly
onto the vanishing sphere $\sqrt{s}\,S^3$, and the involution
acts antipodally.  Passing to the free quotient gives the asserted
retraction onto $\mathbb{RP}^3$ and its rational homology.
The link is the circle bundle with Euler class
$K_{\PP^1\times\PP^1}=-2f_1-2f_2$.  In its homotopy sequence
the image of $\pi_2(\PP^1\times\PP^1)\to\ZZ$ is $2\ZZ$,
so $\pi_1(L)=\ZZ/2$ and $H^1(L;\ZZ)=0$.
The rational circle-bundle sequence gives
$H_2(L;\QQ)=\QQ(f_1-f_2)$.  Lefschetz duality gives
$\dim H_3(M,L;\QQ)=\dim H^3(M;\QQ)=1$; the pair sequence
and $H_2(M;\QQ)=0$ make its boundary map surjective onto the
one-dimensional link group.  It is therefore an isomorphism.
\end{proof}

Smoothness in Lemma~\ref{lem03:quadric} does not assert vanishing of
cotangent cohomology in degree two.  In fact the local calculation gives
$\dim T^1_Q=1$ and $\dim T^2_Q=6$.  The latter value, reproduced in
\texttt{certificates/d19\_quadric\_partial\_resolution.py}, is not needed
in the proof: the explicit equations and full versality show directly
that the obstruction maps vanish.  In particular no complete-intersection
hypothesis on $Q$ is being used.

\subsection{Central fibres and their curve classes}

For the degree-six surface use the marking
\[
 E=\operatorname{Bl}_{p_1,p_2,p_3}\PP^2,
 \qquad h^2=1,\quad e_i^2=-1,\quad
 h\cdot e_i=e_i\cdot e_j=0\ (i\ne j),
\]
with the three points noncollinear.  Its toric boundary, in cyclic
order, is
\begin{equation}\label{eq03:boundary}
 e_1,\quad h-e_1-e_2,\quad e_2,\quad h-e_2-e_3,
 \quad e_3,\quad h-e_1-e_3.
\end{equation}
We retain the root conventions
\begin{equation}\label{eq03:roots}
 \ell=-h+e_1+e_2+e_3,\qquad
 A=-e_1+e_3,\qquad B=e_2-e_3.
\end{equation}
Thus the $A_1$ summand is $\QQ\ell$, and the $A_2$ summand is
$\QQ A\oplus\QQ B$.

\begin{proposition}\label{prop03:localmodels}
There are proper crepant partial resolutions with the following
properties.
\begin{enumerate}
\item A small morphism $W_7\to C_7$ has precisely one singular point,
of type $Q$.  A smooth crepant resolution through $W_7$ differs from
$\operatorname{Tot}(K_{\dP_7})$ by the flop of $h-e_1-e_2$.
Under this comparison the ruling difference of $Q$ is the degree-seven
root, up to the choice of its sign.
\item A morphism $W_{6,1}\to C_6$ is obtained from
$\operatorname{Tot}(K_E)$ by contracting
$C_0=e_1$ and $C_3=h-e_2-e_3$.  It has exactly two ordinary double
points.  The image of $E$ is $\PP^1\times\PP^1$.
\item A morphism $W_{6,2}\to C_6$ is obtained by contracting
$e_1,e_2,e_3$ in $\operatorname{Tot}(K_E)$.  It has exactly three
ordinary double points.  The image of $E$ is $\PP^2$.
\end{enumerate}
The last two morphisms to $C_6$ are divisorial.  With Hodge-normalized
nodal coefficients, the maps into the selected root subspaces are
\begin{equation}\label{eq03:nodaltransfer}
 \lambda\ell=\lambda[C_0]-\lambda[C_3],\qquad
 aA+bB=-a e_1+b e_2+(a-b)e_3.
\end{equation}
In particular the two coefficients over an $A_1$ point are
$(\lambda,-\lambda)$, and the three coefficients over an $A_2$ point
are $(\mu_1,\mu_2,\mu_3)=(-a,b,a-b)$.
\end{proposition}

\begin{proof}
For degree seven use the height-one pentagon
\[
 (0,0),(1,0),(2,1),(1,2),(0,1).
\]
Subdivide it into the triangle with vertices $(0,0),(1,0),(0,1)$
and the quadrilateral with vertices $(1,0),(2,1),(1,2),(0,1)$.
The triangle is unimodular.  Translation of the quadrilateral by
$(-1,-1)$ gives the polygon in Lemma~\ref{lem03:quadric}.
The subdivision preserves the support and introduces no rays;
the induced toric morphism is therefore proper and small.
Every ray has height one, so the morphism is crepant.  The two maximal
cones give one smooth chart and one $Q$ chart, proving the assertion
about its singularities.

Inserting the interior lattice point of the quadrilateral gives a
smooth resolution through this model.  Its triangulation and the star
triangulation of the pentagon differ by one diagonal change.  The
corresponding curve on the original degree-seven surface is
$c=h-e_1-e_2$.  It is a $(-1)$-curve and $K_E\cdot c=-1$; its normal
bundle in $\operatorname{Tot}(K_E)$ is
$\cO(-1)\oplus\cO(-1)$.  The diagonal change is its ordinary flop.
The surface after the flop is the contraction of $c$, namely
$\PP^1\times\PP^1$.  The classes $h-e_1$ and $h-e_2$ have
intersection zero with $c$, square zero, and mutual intersection one.
They descend to its two rulings.  Their difference is $e_2-e_1$,
the degree-seven root.  This also identifies their classes on the
common complement of the flopping curves.  Indeed representatives of
these two pencils can be chosen disjoint from $c$; their differences
give the same class under the homology comparison across the flop.

For degree six use the polygon
\[
 v_0=(1,0),\ v_1=(1,1),\ v_2=(0,1),\ v_3=(-1,0),
 \ v_4=(-1,-1),\ v_5=(0,-1)
\]
and its centre $o=(0,0)$.  The six triangles
$[o,v_i,v_{i+1}]$ are unimodular; the star fan gives
$\operatorname{Tot}(K_E)$ with boundary marking
\eqref{eq03:boundary}.  To contract the boundary curve at $v_i$,
merge its two adjacent triangles into
$[o,v_{i-1},v_i,v_{i+1}]$.  For $i=0,3$, these mergers are disjoint
in their interiors and leave two unimodular triangles.  For $i=0,2,4$
they give three quadrilaterals covering the polygon.
Each merged cell is a unit parallelogram: its four lifted rays have
the relation with coefficients $(1,-1,1,-1)$ in cyclic order, and
three adjacent lifted rays have determinant $\pm1$.  Its affine
toric chart is consequently $xy-zw=0$, with no finite quotient.
No boundary cone changes, and the subdivisions cover the original
cone, proving properness, crepancy, and the claimed singularity counts.

The same contractions can be checked intrinsically.  Every curve
in \eqref{eq03:boundary} has self-intersection and canonical degree
$-1$.  Its normal sequence in the total canonical bundle splits,
since its extension class belongs to $H^1(\PP^1,\cO)=0$.
The selected curves are disjoint, so their contractions are independent
ordinary nodal contractions.  Contracting $e_1,e_2,e_3$ on $E$ gives
$\PP^2$.  Contracting $e_1,h-e_2-e_3$ gives the smooth toric
degree-eight surface with four boundary curves of square zero,
namely $\PP^1\times\PP^1$.  These surfaces are still contracted
by the subsequent morphism to $C_6$, explaining its divisorial nature.

Finally, \eqref{eq03:nodaltransfer} follows by substitution in
\eqref{eq03:roots}.  These are also the deformation-theoretic transfer
maps: inclusion of the nodal exceptional supports in the full
exceptional surface sends a tuple of local classes to the sum of
their curve classes.  Naturality of the connecting maps in local
cohomology and contraction with the same volume form make this map
commute with the Hodge comparison.  In the degree-seven case the
comparison is made on the common regular open set across the flop.
There is thus no independent scalar choice at each newly introduced
node in \eqref{eq03:nodaltransfer}.
\end{proof}

\subsection{A simultaneous modification of the one-dimensional branch}

The two-node model extends over the entire $A_1$ branch.  Here it is
useful to give the extension explicitly, since its central-fibre
description alone would not justify lifting a family.
Altmann's triangle-plus-triangle family is the affine cone over the
Segre embedding of $(\PP^1)^3$, equipped with the function
\begin{equation}\label{eq03:a1function}
 u_0v_0w_0-u_1v_1w_1=s;
\end{equation}
see \cite[Section~8.2(i)]{Altmann97}.  The products in this equation
are tensor coordinates on the cone, rather than functions on the
product of projective spaces.

Retain the first two projective factors.  The resulting incidence
space is
\[
 V=\operatorname{Tot}_{\PP^1\times\PP^1}
       \bigl(\cO(-1,-1)\oplus\cO(-1,-1)\bigr).
\]
Its morphism to the tensor cone sends a pair of fibre vectors to the
tensor $u\otimes v\otimes w$.  It is proper, since the incidence
description is closed in the product of the cone with
$\PP^1\times\PP^1$.  It is an isomorphism away from the vertex:
a nonzero decomposable tensor determines its first two factor lines
uniquely.  On $V$, equation \eqref{eq03:a1function} defines a morphism
to the $s$-line and hence a simultaneous modification of that family.

There are exactly two critical points in the central fibre.  To see
this, write the fibre coordinates locally as $A,B$.  Differentiating
with respect to them forces $u_0v_0=u_1v_1=0$.  The two possible
base points are $([1:0],[0:1])$ and $([0:1],[1:0])$.  On their
respective affine charts the equations are exactly
\begin{equation}\label{eq03:a1charts}
 \beta A-\alpha B=s,
 \qquad \gamma C-\delta D=s.
\end{equation}
Their differentials vanish only at the origin of each four-dimensional
chart, and their quadratic forms are nondegenerate.  On the other
two projective charts one coefficient of a fibre coordinate is a unit,
so the map is smooth there.  These equations also prove flatness:
near a critical point the morphism is the standard nodal deformation,
and elsewhere it is smooth.  Equivalently, the local rings are
torsion-free over the regular one-dimensional base.

The zero section belongs only to $s=0$.  Thus the modification is an
isomorphism over every $s\ne0$ and away from the vertex of the central
fibre.  Its central fibre is $W_{6,1}$: resolving the two nodes on
the side which blows up the corresponding points of the zero section
gives the marked degree-six surface and the two curves $C_0,C_3$.
This can equally be read from the first degree-six subdivision in
Proposition~\ref{prop03:localmodels}.

The two displayed equations have the same algebraic parameter $s$.
Their Hodge-normalized parameters have opposite signs when their
exceptional curves are oriented by $[C_0],[C_3]$.  More precisely,
the image is $(\lambda,-\lambda)$ with $\lambda=c_1s$ to first
order for a nonzero constant $c_1$.  This follows either from the
local-cohomology comparison above or by pairing with $E$:
$E\cdot C_0=E\cdot C_3=-1$, while the branch class lies in
$K_E^\perp$.  Neither nodal derivative vanishes, by
\eqref{eq03:a1charts}.  The normalized branch coordinate is chosen
using this $c_1$, rather than by equating two independently chosen
residue normalizations.

\subsection{The two-dimensional branch and its discriminant}

\begin{proposition}\label{prop03:a2family}
The $A_2$ branch of $C_6$ admits a proper simultaneous partial
resolution whose central fibre is $W_{6,2}$.  The morphism is an
isomorphism over every nonzero point of the two-dimensional base,
including its discriminant.  In linear coordinates $(a,b)$ compatible
with the Hodge normalization, the three nodal tangent parameters are
\[
 (\mu_1,\mu_2,\mu_3)=(-a,b,a-b).
\]
The discriminant is the union of the three lines
$ab(a-b)=0$.  At a nonzero point of any one of these lines the fibre
has precisely one ordinary double point; off these lines it is smooth.
\end{proposition}

\begin{proof}
We first work in algebraic parameters before imposing the Hodge
normalization.  Write $D=\CC^3/\CC(1,1,1)$ and represent its points
by $(d_1,d_2,d_3)$.  The branch family is
\begin{equation}\label{eq03:matrixfamily}
 \mathcal M_d=\{M:\ \rk M\leq1,\quad M_{ii}=z+d_i\ (1\leq i\leq3)\}.
\end{equation}
Common translation of the $d_i$ is absorbed in $z$.  This is the
actual restricted versal family, not just a first-order model.
For an explicit comparison with \cite[Section~5.6]{Altmann97}, take
edge parameters $(a,b,c,a,b,c)$ and the matrix
\[
 \begin{pmatrix}
 z+c&z_1&z_2\\ z_4&z+b&z_3\\ z_5&z_6&z+a
 \end{pmatrix}.
\]
Its nine $2\times2$ minors are the six equations
$z_it_i-z_{i-1}z_{i+1}=0$ (cyclic indices), with
$(t_1,\ldots,t_6)=(z+a,z+b,z+c,z+a,z+b,z+c)$,
and the three equations
$t_2t_3-z_1z_4=t_3t_4-z_2z_5=t_1t_2-z_3z_6=0$.
They generate exactly the branch ideal.  This also follows from
\cite[Section~8.2(ii)]{Altmann97}, which describes its total space
as the cone over $\PP^2\times\PP^2$ with the two diagonal-difference
functions.

Resolve that total space by retaining its column line:
\begin{equation}\label{eq03:incidence}
 \widetilde{\mathcal M}
 =\{(M,L):\operatorname{im}M\subseteq L\}
 =\operatorname{Tot}_{\PP^2}(\cO(-1)^{\oplus3})
 \longrightarrow\{M:\rk M\leq1\}.
\end{equation}
Write $M=xy^{\mathsf T}$, where $x\ne0$ and
$(x,y)\sim(tx,t^{-1}y)$.  The incidence relation proves properness.
A nonzero rank-one matrix has a unique column line, and on an open
set $M_{k\ell}\ne0$ this line is represented by its $\ell$th column.
These formulas give regular inverses, proving that
\eqref{eq03:incidence} is an isomorphism away from the zero matrix.

Compose \eqref{eq03:incidence} with the diagonal-difference map to
$D$.  On $x_k=1$, let $i,j$ be the two other indices.  Elimination of
$y_k,z$ gives the single equation
\begin{equation}\label{eq03:a2chart}
 x_i y_i-x_j y_j=d_i-d_j,
 \qquad y_k=x_i y_i-d_i+d_k.
\end{equation}
The two independent base coordinates can be taken as $d_i-d_j$
and $d_k-d_i$.  Thus each chart is the product of the standard nodal
deformation with a smooth one-dimensional base.  This proves
flatness over the full plane, including its origin and discriminant.
In the central fibre its only singular point in this chart has all
four coordinates zero.  The three resulting points lie over the
three distinct coordinate points of $\PP^2$ and give exactly three
nodes.  The zero section is the exceptional $\PP^2$.  Resolving the
nodes by blowing up their points on this surface recovers
$\operatorname{Bl}_3\PP^2$ and the curves $e_1,e_2,e_3$.
The star fan is the second subdivision of
Proposition~\ref{prop03:localmodels}; in particular the central
morphism is crepant and is $W_{6,2}\to C_6$.

The zero matrix has all diagonal differences zero.  It therefore
occurs only over $0\in D$.  The isomorphism away from the zero
matrix consequently restricts to an isomorphism on every fibre with
$d\ne0$, with no exception along a discriminant line.  The chart
equations show that the three singularity parameters are the cyclic
differences
\[
 d_2-d_3,\quad d_3-d_1,\quad d_1-d_2.
\]
Their sum is zero.  A nonzero base point lies on at most one of their
zero lines, and the corresponding chart has one node by its
nondegenerate quadratic equation.  All remaining charts are smooth.
This proves both the discriminant statement and the assertion about
its noncentral fibres.

To identify parameters with the fixed marking, label the three
coordinate points, and choose signs in the nodal equations, so that
their differences in the order $e_1,e_2,e_3$ are
$d_3-d_1,d_2-d_3,d_1-d_2$.  Taking
$(d_1,d_2,d_3)=(a,b,0)$ gives $(-a,b,a-b)$ algebraically.
The determinant-twisted equivariance in
Lemma~\ref{lem:branchhodge} identifies the two-dimensional parameter
representation with the $A_2$ root representation, uniquely up to a
nonzero scalar: rotation and reflection act as the usual irreducible
two-dimensional representation of the triangle symmetry group.
Thus, after matching this marking, the Hodge comparison multiplies
these three algebraic differences by one common nonzero scalar.
Rescale $a,b$ simultaneously by that scalar.  Naturality under
inclusion of supports now gives exactly
\[
 aA+bB=-a e_1+b e_2+(a-b)e_3
\]
with the prescribed global volume form.  The linear rescaling leaves
the three discriminant lines unchanged.  This proves the asserted
normalized parameter map.
\end{proof}

The constants relating algebraic parameters to Hodge-normalized
parameters may differ between singular points, and between the
$A_1$ and $A_2$ summands at one degree-six point.  They are fixed by
the one global volume form, not by unrelated choices of local
coordinates.  Equations~\eqref{eq03:nodaltransfer} are identities
after this normalization.  Their vanishing conditions are independent
of the remaining common nonzero factors.

The simultaneous constructions also distinguish a nonzero tangent
on a smoothing component from a tangent transverse to its
discriminant.  Every nonzero $A_1$ parameter smooths the germ, whereas
a nonzero vector in the $A_2$ plane may leave one of the three nodes
unsmoothed.  Because the simultaneous modification is an isomorphism
on every noncentral fibre, this distinction can be applied to an
arbitrary analytic classifying arc, including an arc tangent to a
discriminant line to high order.  It is this property that permits
the necessity argument in Section~\ref{sec04:smoothing}.
